% Section 5.2, from lectures/L05/README.md: the goals, the evaluation questions, and the next
% lecture.

\section{Review}
\label{c5:sec:review}

\needspace{6\baselineskip}
\subsection*{What you should be able to do now}
After this chapter, you should:
\begin{itemize}
  \item Explain why a busy-wait delay (\code{_delay_ms}) is unsuitable once a program needs to do
    more than one thing at a time.
  \item Be able to compute a timer's tick period from the CPU clock and a chosen prescaler.
  \item Be able to work out how many overflow interrupts are needed for a given delay, for an
    8-bit timer.
  \item Understand the practical difference between Normal Mode and CTC Mode, and when you'd reach
    for each.
  \item Know why a variable shared between an ISR and \code{main()} needs to be declared
    \code{volatile}.
\end{itemize}

\needspace{6\baselineskip}
\subsection*{Questions to test yourself}
You should be able to answer:
\begin{enumerate}
  \item At 16\,MHz with a \textdiv256 prescaler, what's the resulting tick period?
  \item What's the difference between what triggers \code{TIMERn_OVF_vect} and what triggers
    \code{TIMERn_COMPA_vect}?
  \item Why is \code{volatile} optional on \code{ovf_count} in the Timer 0 example, and when does a
    variable shared with an ISR genuinely need it?
\end{enumerate}

\needspace{6\baselineskip}
\subsection*{Looking ahead}
Chapter~\ref{c6:ch} is about serial communication:
\begin{itemize}
  \item Serial communication (UART): transmitting data to a PC instead of only driving LEDs.
  \item Combining timers with UART to report readings periodically.
\end{itemize}
